This chapter relates BlockQuiz to the literature it draws on.
It reviews the research on block-based programming and learning, surveys related
platforms and lays out the educational theories that informed the design.

\section{Block-Based Programming and Learning}
Block-based programming environments represent program structure through shaped, snappable pieces
instead of textual syntax. The shapes encode which constructs fit together.
Loop blocks accept statements, expressions return values into matching sockets and the editor refuses
combinations that the underlying language would reject as a parse error.
Maloney et al. describe the Scratch environment around several specific design choices.
liveness with no compile or run mode, tinkerability through clickable blocks, visible execution, a deliberate absence of 
error messages
because block shapes prevent syntactically meaningless connections, concrete data through on-stage variable monitors 
and a minimised command set
\cite{maloney2010}. Resnick and colleagues frame the wider goals as three design principles for Scratch. Making it
\texttt{more tinkerable, more meaningful and more social} than other programming environments \cite{resnick2009}.

Weintrop and Wilensky followed high-school students for ten weeks in an introductory programming course that began in Snap! 
\cite{snap2026} and transitioned to Java and they asked students what they themselves thought of the two modalities 
\cite{weintrop2015}. Their analysis identified four reasons students gave for block-based programming
being easier. The natural-language labels on blocks, the visual shape and layout of the blocks, the ease of composing 
programs through drag and drop and the use of blocks as memory aids since the toolbox lets a student browse the available 
commands. About 58 percent of the 84 mid-study responses identified blocks as easier than the text-based equivalent.
Students also named drawbacks. Blocks felt less authentic, less powerful and less suited to trial-and-error exploration 
than text-based programming.
The paper is qualitative. Its design implication is not that blocks are simply better.
The authors argue that the question for a learning environment is how to keep the introductory benefits while staying 
honest about where blocks get in the way.

Several empirical studies have looked at what learners actually take away from short Scratch interventions. 
Kalelioglu and Gülbahar ran a five-week, one-hour-per-week unit
with 49 Turkish fifth-grade students and reported mixed results.
Their quantitative measure of problem-solving skill did not show a significant change after the unit, while the 
qualitative responses from the children
themselves were positive about Scratch and about programming as an activity \cite{kalelioglu2014}.
Saez-Lopez and colleagues ran a longer two-year case study with 107 primary students across five Spanish schools and 
reported
significant improvements in computational concepts such as sequences, loops, parallelism and events, together with
positive attitudes toward learning with the platform \cite{saezLopez2016}. Meerbaum-Salant et al. developed 
Scratch-based learning materials
for middle-school students and evaluated them in normal classes at two schools \cite{meerbaum2013}.
Combining written tests with classroom observation, the study reported that students could successfully learn key 
computer-science concepts in Scratch but also
ran into recurring problems around initialisation, variables and concurrency.
Across these studies the recurring pattern is that the measured learning gains are more modest than the positive responses 
from learners would suggest.

\subsection{Computational Thinking as a Learning Goal}
The wider research frame for these studies is \texttt{computational thinking}. Wing framed the term as a fundamental
skill for everyone, characterised as the thought process of formulating a problem and expressing its solution so that an
information-processing agent can carry it out \cite{wing2006}. Grover and Pea reviewed
the six years of research after Wing's article and reported on the state of the field, the gaps in the research and
priorities for future inquiry into K-12 (primary and secondary school) computational thinking \cite{grover2013}.
Brennan et al. proposed an operational framework for studying computational thinking through Scratch \cite{brennan2012}.
They divide the construct into three layers. Computational concepts cover sequences, loops, parallelism, events, conditionals, operators
and data. Computational practices cover the way learners build a program over time. Being incremental and iterative, testing and debugging, reusing and
remixing and abstracting and modularising.
Computational perspectives describe how learners position themselves in relation to the medium, including expressing themselves, connecting with others
and questioning the technology around them.
The Brennan and Resnick framework matters for the present work in two ways. First,
it argues that debugging is not a side activity but a core computational practice.
A learning environment that shows only a pass or fail signal forfeits a large part of the learning opportunity.
Second, the framework treats the artefact and the process around it as separate sources of evidence.
A platform that records what a learner did, in addition to whether they succeeded, produces richer data than one that
records outcomes alone.


\subsection{Educational Block-Based Environments}


\subsection*{Scratch}
Scratch \cite{scratch2026}, developed at the MIT Media Lab, is one of the most influential systems in this area \cite{resnick2009}.
Its educational strength lies in the way it connects programming
to storytelling, games, animation and personally meaningful projects.
The Scratch website hosts a large community in which learners share and remix projects, which Brennan and Resnick drew on in their work on computational-thinking assessment
through project portfolios, artefact-based interviews and design scenarios \cite{brennan2012}. The same openness that makes Scratch pedagogically rich makes it
difficult to use for short tasks with clear expected outcomes.
Scratch does not ship with a deterministic auto-grading mechanism.
Weintrop and Wilensky note that students also perceive block environments as less suited to text-style trial-and-error than text-based programming,
an observation that matters for any short-task setting in which a quick second attempt is part of the intended learning loop \cite{weintrop2015}.

\subsection*{Microsoft MakeCode}
Microsoft MakeCode \cite{makecode2026} extends the idea of block-based learning toward hardware and project-based exploration.
It offers polished editors, simulators and a transition path from blocks to JavaScript or Python.
This makes it particularly useful in contexts where physical computing or the blocks-to-text progression is part of the learning goal.
However, its overall focus differs from a platform centered on short, tightly constrained exercises.
In comparison to the present work, MakeCode is less concerned with teacher-authored assessment tasks and more concerned with broad exploratory creation in hardware-oriented settings.


\subsection*{Google Blockly}
Google Blockly is best understood not as a complete learning platform but as a library for building visual programming interfaces \cite{blockly}. This makes it particularly relevant for
custom educational systems. Blockly provides a mature rendering and connection model for blocks, configurable toolboxes, internationalization support, and code generators for multiple
target languages. Because it is a library rather than a fixed educational product, it enables the construction of environments in which the available blocks, the surrounding
user interface and the grading logic can be aligned with specific pedagogical goals. For BlockQuiz, Blockly therefore serves as a technical foundation rather than as an end-user product.

\subsection*{Code.org}
Code.org \cite{codeorg2026} demonstrates how strongly structured learning paths can support school teaching. Its courses combine guided exercises,
immediate feedback, and clear curricular progression, which makes it an important reference point for guided beginner instruction.
In contrast to more open creative environments, Code.org is designed around sequenced lessons and classroom management. This makes it particularly relevant as a
comparison for BlockQuiz, since both systems emphasize guided tasks over unrestricted project work. However, Code.org is primarily delivered as a centralized
online platform and is less focused on self-hosted, privacy-conscious deployment in local school environments.

\subsection*{Open Roberta}
Open Roberta \cite{roperta} is especially relevant from a European and school-oriented perspective. Developed as a browser-based platform for educational robotics,
it was designed to reduce technical barriers for teachers and learners by avoiding complex local installation and configuration steps.
Its visual programming model, accessibility for beginners, and compatibility with classroom contexts make it a valuable comparison point for BlockQuiz.
At the same time, Open Roberta is centered more strongly on programming robots and boards, whereas BlockQuiz focuses on constrained programming exercises with
integrated auto-grading and teacher-authored task sequences.

\subsection*{LEGO Mindstorms}
LEGO Mindstorms \cite{lego} is one of the best-known robotics platforms in introductory programming education and has played an important role in bringing
programming into schools through hands-on, tangible activities. Its strength lies in the connection between software and physical devices,
which can be highly motivating for learners. In educational terms, it shares with Open Roberta a strong emphasis on robotics, experimentation,
and visible program behavior. However, this hardware-centered approach also distinguishes it from BlockQuiz, whose main purpose is not
robot control but the delivery of guided, browser-based programming exercises that can be deployed and assessed without requiring dedicated physical equipment.

\subsection{Comparison Summary}
Taken together, these platforms show that there is no single dominant design for block-based programming education. Scratch foregrounds creativity, while MakeCode, Open Roberta and LEGO Mindstorms bridge visual and textual programming
in hardware-oriented contexts. Code.org offers strongly scaffolded instruction and Blockly acts as an enabling framework.
BlockQuiz positions itself with a different focus where constrained exercises, teacher authoring, automatic assessment and privacy-conscious deployment need to coexist.

The matrix in Table~\ref{tab:platform-comparison} shows that BlockQuiz does not compete with these platforms on creative freedom. It instead occupies a specific intersection that is otherwise empty.
Per-exercise toolbox constraints under a GUI authoring workflow, automatic grading, structured per-attempt analytics that go beyond completion counts and a deployment model that is
genuinely self-hostable inside school infrastructure rather than only theoretically open source.
Scratch and Blockly Games optimize for open creation rather than guided practice. MakeCode targets hardware tutorials. None of them combine teacher-side authoring of constrained tasks with on-prem deployment
and learner-process analytics.

\begin{table}[H]
  \centering
  \caption{Comparison of Block-Based Programming Platforms}
  \label{tab:platform-comparison}
  \begin{tabular}{|l|c|c|c|c|}
    \hline
    \textbf{Platform}  & \textbf{Auto-grading} & \textbf{Teacher CMS} & \textbf{Self-hostable} & \textbf{Open Source} \\
    \hline
    Scratch            & No                    & No                   & No                     & Yes                  \\
    MakeCode           & Partial               & No                   & No                     & Yes                  \\
    Code.org           & Yes                   & Limited              & No                     & Partial              \\
    Blockly            & N/A (library)         & N/A                  & Yes                    & Yes                  \\
    \textbf{BlockQuiz} & Yes                   & Yes                  & Yes                    & Yes                  \\
    \hline
  \end{tabular}
\end{table}

\section{Automated Assessment in Programming Education}
\label{sec:auto-grading}

Automated assessment has a long history in programming education and
provides immediate feedback at scale. Early work in the field dates
from the 1960s, with classroom systems described by Hollingsworth
\cite{hollingsworth1960}, Naur \cite{naur1964}, and Forsythe and Wirth
\cite{forsythe1965}. Reviews of more recent decades document the
continued development of the area
\cite{alaMutka2005, douce2005, ihantola2010, messer-neil}.
Ihantola and colleagues survey systems from 2006 to 2010 and discuss
the pedagogical and technical features they support
\cite{ihantola2010}. Across these reviews, the recurring observation
is that most tools assess correctness through dynamic test execution.


\subsection{Test-Based Grading}
The most common approach executes student code against a set of test cases~\cite{douce2005, alaMutka2005},
comparing actual output to expected outputs. This is well suited to input/output exercises in which correctness can be defined through a clear mapping from input to output.
In practice, such grading must deal with issues such as visible versus hidden tests, normalization of whitespace or upper/lower case and the distinction between purely summative checking
and formative feedback.
The systematic review by Messer et al. shows that this form of grading remains the dominant approach in automated programming assessment because it
is reliable for constrained tasks and scales well to large cohorts \cite{messer-neil}.

\subsection{State-Based Grading}
Not all beginner tasks produce meaningful textual output. In graphical or simulation-based exercises, correctness is often better captured through the final state of the system
or through the sequence of actions that produced the state.
Turtle graphics are a good example: a learner may arrive at the same picture through multiple valid block sequences, which means that grading can
be based either on the final geometric state or a more prescriptive comparison of executed commands \cite{malmi2004}. For this reason, BlockQuiz combines simple command-based and state-based
grading strategies, depending on the intended learning objective.

\subsection{Benefits of Immediate Feedback}
Formative-assessment research treats feedback as one of the most influential factors on classroom learning. 
Black and Wiliam reviewed the empirical evidence on classroom assessment and concluded that strengthening the feedback 
students receive about their learning yields substantial learning gains \cite{blackWiliam1998}. Hattie and Timperley 
summarised feedback research at four levels (tasks, process, self-regulation, self) and reported that feedback aimed at the 
task and at the learner's process is among the most powerful influences on learning \cite{hattie2007}. Shute synthesised the literature 
on formative feedback and reported that the most effective timing depends on task complexity and learner ability. 
Immediate feedback tends to be more efficient for low-achievement learners and for simple or procedural tasks, 
whereas delayed feedback can favour high-achievement learners and transfer to new problems \cite{shute2008}. For 
the BlockQuiz target setting, beginner-level learners working on short, procedural exercises, Shute's analysis points 
toward the immediate end of the timing range. 

\section{Learning Theory Foundations}
The design of BlockQuiz is informed by several well-established educational theories.
These theories do not determine the implementation directly, but they provide a conceptual framework for understanding why constrained block-based
tasks, progressive hints and visual feedback are appropriate design choices.

\subsection{Constructionism}
Seymour Papert's constructionism argues that learning becomes especially powerful when students actively construct artifacts
that are meaningful to them \cite{papert1980}. Programming can therefore be understood not merely as a way of issuing instructions to a machine, but as a
medium through which learners explore ideas, observe consequences and build mental models.
Although BlockQuiz focuses on shorter and more guided tasks than open creative environments, it still adopts the constructionist insight that programming is most
engaging when learners can see the effects of their decisions directly. Turtle graphics, in particular, make this visible by connecting code to movement and shape.

\subsection{Cognitive Load Theory}
Sweller's Cognitive Load Theory distinguishes between the inherent difficulty of a task, extraneous load introduced by presentation or interface design and
beneficial load associated with meaningful learning activity \cite{sweller1988}.
One of the main strengths of block-based programming is that it reduces extraneous load by eliminating syntax errors and making valid composition visible.
BlockQuiz extends this idea by constraining the toolbox on a per-exercise basis.
Instead of exposing a large library of irrelevant blocks, the system presents only the constructs that are needed for a given task, which
supports focused learning and reduces distraction.


\subsection{Zone of Proximal Development}
Vygotsky's Zone of Proximal Development describes the range in which learners can complete a task with guidance but not yet entirely on their own \cite{Vygotsky1980, Shabani}.
This idea is relevant for platforms that provide progressive support rather than all-or-nothing assessment. In BlockQuiz, the use of staged hints, starter blocks and ordered
task progression is intended to keep exercises within that zone. The platform does not merely judge whether an answer is correct, it is also
designed to support learners while they move toward a correct solution.



\subsection{Piaget's Stages of Cognitive Development}
The 8 to 12 age range sits mostly inside Piaget's concrete operational stage (roughly ages 7 to 11) and overlaps the start 
of the formal operational stage at the upper end of the band \cite{PiagetInhelder1969}. Learners 
in the concrete operational stage reason confidently about concrete objects and relationships but find purely abstract symbolic reasoning harder. 
Block-based programming maps well onto this stage. The blocks are concrete objects on the screen. 
Their shapes and colours stand in for the relationships that text-based syntax encodes only in symbols. 
The turtle and grid-robot exercises take this further. 
The program acts on a visible object whose position, orientation and collected items the learner observes directly. 
BlockQuiz uses this mapping deliberately. The visual outcome of a program is the channel through which learners 
check their own work, before they consult any written feedback. 

\subsection{Summary}
Taken together, these theories shape four design moves in BlockQuiz. Per-exercise toolbox constraints 
follow Cognitive Load Theory and the empirical findings on block environments. 
Progressive hints follow Vygotsky and Bruner. Visible, immediate outcomes follow constructionism and Piaget. The next chapter derives 
the requirements that turn these design moves into concrete platform features. 
